% Section 5: Per-Irrep Analysis and Domain-Dependent Activation
% Draft by Lyra -- Sections 5.1 and 5.2
% For discussion with Clio before integration into main paper.
%
% Assumes LNCS or similar style with amsmath, amssymb loaded.
% Clio will handle Section 5.3 (fiber rotation mechanism).

\section{Representation-Theoretic Analysis of the Cactus Action}
\label{sec:irrep-analysis}

The headline result of Section~4 established a strict ordering of
commutator rank by interval arrangement, uniform across all non-trivial
irreducible representations of~$S_n$.  We now ask two deeper questions.
First, how does this ordering decompose when we resolve it per
irreducible representation? (Section~\ref{sec:per-irrep-rank}.)
Second, why does the experimental sign of the arrangement effect
\emph{reverse} across task domains?
(Section~\ref{sec:domain-dependent-activation}.)

A guiding intuition throughout this section is the duality between
\emph{sheaves} and \emph{copresheaves} on communication graphs.
Sheaves enforce consistency: a sheaf assigns data to nodes subject to
agreement conditions on overlaps, formalizing the coordination channel
in multi-agent systems.  Copresheaves, by contrast, enable directed
information flow \emph{without} consistency constraints, formalizing
the exploration channel~\cite{hajij2025copresheaf}.
The signed laxator, as we shall see, can be understood as a natural
transformation mediating between these two perspectives: its sign
determines whether the system is driven toward sheaf-like consensus or
copresheaf-like divergence.

%% ================================================================
\subsection{Per-Irrep Rank Ordering}
\label{sec:per-irrep-rank}

\subsubsection{Setup and notation.}

Recall from Section~3 that each cactus generator
$s_{p,q} \in J_n$ acts on the Specht module~$V_\lambda$ of~$S_n$ via
the longest element $w_0^{[p,q]} \in S_{[p,q]}$, which reverses the
interval~$[p,q]$.  For each pair of generators
$(s_{p_1,q_1},\; s_{p_2,q_2})$, the \emph{commutator rank} is
\[
  c_\lambda(s_{p_1,q_1},\, s_{p_2,q_2})
  \;=\;
  \operatorname{rank}\!\bigl(
    \rho_\lambda(s_{p_1,q_1})\,\rho_\lambda(s_{p_2,q_2})
    \;-\;
    \rho_\lambda(s_{p_2,q_2})\,\rho_\lambda(s_{p_1,q_1})
  \bigr),
\]
and the pair is classified by the relative arrangement of
$[p_1,q_1]$ and $[p_2,q_2]$ into one of four types:
\emph{Separated} ($q_1 < p_2$, disjoint),
\emph{Overlapping} (partial overlap, neither contains the other),
\emph{Nested} (one strictly contains the other), or
\emph{Adjacent} (intervals share exactly one endpoint).

\subsubsection{The strict ordering at $n = 5$.}

Table~\ref{tab:n5-commutator-rank} presents the mean commutator rank
by arrangement type for every non-trivial irrep of~$S_5$.  The
computation covers all~$\binom{10}{2} = 45$ pairs of the ten cactus
generators of~$J_5$, evaluated on each of the five non-trivial Specht
modules (390 rows total; see Appendix~A for the complete table).

\begin{table}[t]
\centering
\caption{Mean commutator rank by arrangement type, per non-trivial
irrep of~$S_5$.  Adjacent generators achieve the maximum commutator
rank in every case; Separated generators always commute.}
\label{tab:n5-commutator-rank}
\smallskip
\begin{tabular}{lcrrrr}
\hline
$\lambda$ & $\dim V_\lambda$ & Separated & Overlapping & Nested & Adjacent \\
\hline
$[4,1]$     & 4 & 0.00 & 1.60 & 1.92 & 2.00 \\
$[3,2]$     & 5 & 0.00 & 2.40 & 2.72 & 4.00 \\
$[3,1,1]$   & 6 & 0.00 & 2.40 & 3.52 & 4.00 \\
$[2,2,1]$   & 5 & 0.00 & 2.40 & 2.72 & 4.00 \\
$[2,1,1,1]$ & 4 & 0.00 & 1.60 & 1.92 & 2.00 \\
\hline
\end{tabular}
\end{table}

\noindent
The ordering
\begin{equation}
\label{eq:strict-ordering}
  \text{Separated} \;<\; \text{Overlapping} \;<\; \text{Nested} \;<\; \text{Adjacent}
\end{equation}
is \emph{strict for every non-trivial irrep} at $n = 5$.  Moreover,
the Adjacent commutator rank is perfectly consistent: every Adjacent
pair achieves the same value within each irrep (rank~$2$ for the
four-dimensional irreps $[4,1]$ and $[2,1,1,1]$, and rank~$4$ for the
five- and six-dimensional irreps $[3,2]$, $[2,2,1]$, and $[3,1,1]$).
Separated generators commute exactly ($c_\lambda = 0$ for all pairs,
all irreps).

\begin{remark}[Adjacent as maximal non-commutativity]
For the five-dimensional irreps $[3,2]$ and $[2,2,1]$, the Adjacent
commutator rank equals~$4$, which is the maximum possible for a
$5 \times 5$ commutator matrix (a commutator $[A,B] = AB - BA$ has
trace zero, so its rank is at most $\dim V_\lambda - 1$ when
$\dim V_\lambda$ is odd, and at most $\dim V_\lambda$ when even; here
$\dim = 5$ gives rank $\leq 4$).  Adjacent generators are thus
\emph{maximally non-commuting} on these irreps.
\end{remark}

\subsubsection{The $n = 4$ degeneracy.}

At $n = 4$, the strict ordering~\eqref{eq:strict-ordering} partially
collapses: Separated and Overlapping generators have identical mean
commutator rank (both zero) for all non-trivial irreps.  This
degeneracy arises because $n = 4$ admits only a single Overlapping pair
--- the intervals $[1,3]$ and $[2,4]$ --- and the corresponding
permutations $w_0^{[1,3]}$ and $w_0^{[2,4]}$ happen to commute in
$S_4$.  (Concretely, these are $(1\;3)$ and $(2\;4)$, which are
disjoint transpositions composed with a fixed point, so they commute on
the nose.)

\begin{table}[t]
\centering
\caption{Mean commutator rank at $n = 4$.  The Separated--Overlapping
degeneracy is a small-$n$ artifact; see text.}
\label{tab:n4-commutator-rank}
\smallskip
\begin{tabular}{lcrrrr}
\hline
$\lambda$ & $\dim V_\lambda$ & Separated & Overlapping & Nested & Adjacent \\
\hline
$[3,1]$   & 3 & 0.00 & 0.00 & 1.78 & 2.00 \\
$[2,2]$   & 2 & 0.00 & 0.00 & 0.89 & 2.00 \\
$[2,1,1]$ & 3 & 0.00 & 0.00 & 1.78 & 2.00 \\
\hline
\end{tabular}
\end{table}

\noindent
The degeneracy breaks cleanly at $n = 5$, where five independent
Overlapping pairs exist and all achieve non-zero mean commutator rank.
This is significant: it confirms that the four-class ordering is a
structural feature of the cactus action, not a coincidence at a single
value of~$n$, and that the $n = 4$ collapse is the artefact rather than
the rule.

\subsubsection{Universality across irreps.}

The key observation is that the ordering~\eqref{eq:strict-ordering}
does not depend on the choice of irrep~$\lambda$.  While the absolute
values of $c_\lambda$ vary (e.g., $[3,1,1]$ achieves a Nested mean of
$3.52$ versus $1.92$ for $[4,1]$), the \emph{relative ranking} of
arrangement types is identical for every non-trivial $\lambda$ at a
given~$n$.  This universality has a concrete interpretation:

\begin{quote}
\emph{For any representation through which the cactus group acts
non-trivially, the degree to which two generators fail to commute is
determined by the geometric relationship of their intervals, not by the
algebraic details of the representation.}
\end{quote}

This is precisely the structure needed for the signed laxator: the
arrangement classification provides an invariant that is simultaneously
algebraic (via the cactus group) and geometric (via interval nesting),
and this invariant robustly predicts the strength of non-abelian
interference across all representation channels.

In the language of sheaf--copresheaf duality, the universality of
the ordering across irreps reflects a structural fact: the
arrangement type controls not only the \emph{degree} of
non-commutativity but also the balance between consistency
enforcement (the sheaf channel) and unconstrained information flow
(the copresheaf channel).  Separated generators, which commute
exactly, preserve the sheaf's consistency conditions; Adjacent
generators, which are maximally non-commuting, break these conditions
maximally and drive the system toward the copresheaf regime.  The
signed laxator~$\phi_G$ sits at the interface, with its sign encoding
which regime dominates for a given topology and task.

%% ================================================================
\subsection{Domain-Dependent Irrep Activation}
\label{sec:domain-dependent-activation}

\subsubsection{The maze sign-flip.}

The universal ordering~\eqref{eq:strict-ordering} predicts that
Adjacent arrangements should always produce the \emph{strongest}
interference effect, and Separated the weakest.  In our experimental
data on NK landscapes and sudoku, this prediction holds: the
interference experiment (Table~4 of the companion data set%
\footnote{Available at
\texttt{github.com/lyra-claude/signed-laxator}, directory
\texttt{experiments/}.}) measures $\eta^2$ for the arrangement effect
at constant $\beta_1 = 2$ across three cycle configurations
(Adjacent, Separated, Nested), and the saturation ordering on sudoku
is
\[
  \text{Separated}~(0.385) \;>\;
  \text{Adjacent}~(0.165) \;>\;
  \text{Nested}~(0.053),
\]
which matches the cactus-rank prediction when we interpret higher
commutator rank as stronger interference and thus higher saturation
ceiling.

% NOTE: The sudoku ordering is Separated > Adjacent > Nested for the
% saturation *ceiling*, which corresponds to the cactus prediction
% that Separated pairs commute (additive, no cancellation) while
% Adjacent pairs have maximal non-commutativity (destructive
% interference lowers the ceiling).  This needs careful framing ---
% see discussion below.

However, on maze tasks the ordering \emph{reverses}:
\[
  \text{Adjacent}~(0.028) \;>\;
  \text{Nested}~(0.027) \;>\;
  \text{Separated}~(0.022).
\]
The effect sizes are small (near the noise floor, with $\eta^2 \approx
0.01$--$0.03$ and diversity locked near $0.989$ across all conditions),
but the reversal is consistent across generations.  On OneMax (NK0),
a partial match obtains:
\[
  \text{Separated}~(0.074) \;>\;
  \text{Nested}~(0.037) \;>\;
  \text{Adjacent}~(0.004).
\]

\subsubsection{Interpretation: the active irrep depends on the task.}

We propose that the resolution lies not in the cactus algebra itself
but in which irreducible representation of the cactus group is
\emph{activated} by a given task domain.  The argument runs as
follows.

The cactus group $J_n$ acts simultaneously on every irrep $V_\lambda$
of~$S_n$.  But the evolutionary dynamics on a particular fitness
landscape do not excite all irreps equally.  Different task domains
impose different symmetry-breaking patterns on the population, and
these patterns project preferentially onto different irreps.

\begin{definition}[Active irrep]
\label{def:active-irrep}
Let $\mathcal{D}$ be a task domain and $\pi_\lambda$ the projection
onto the $\lambda$-isotypic component of the population state space.
We say $V_\lambda$ is \emph{$\mathcal{D}$-active} if
$\|\pi_\lambda(\delta_{\mathcal{D}})\| > \epsilon$ at the relevant
timescale, where $\delta_{\mathcal{D}}$ is the diversity trajectory
difference vector across arrangements and $\epsilon$ is a
significance threshold.
\end{definition}

The per-irrep commutator ranks in Table~\ref{tab:n5-commutator-rank}
are universal in their \emph{ordering} but not in their \emph{absolute
values}.  The spread between arrangement types differs:
$[3,1,1]$ has Nested mean~$3.52$ versus Adjacent~$4.00$ (ratio
$0.88$), while $[4,1]$ has Nested~$1.92$ versus Adjacent~$2.00$
(ratio~$0.96$).  If different task domains activate different
irreps, the \emph{observable} arrangement effect is the
weighted sum
\begin{equation}
\label{eq:weighted-arrangement}
  \eta^2_{\mathrm{obs}}(\text{arrangement})
  \;=\;
  \sum_\lambda
    w_\lambda(\mathcal{D})\;
    f\bigl(c_\lambda(\text{arrangement})\bigr),
\end{equation}
where $w_\lambda(\mathcal{D})$ encodes the domain-dependent activation
weights and $f$ is a monotone function relating commutator rank to
diversity effect.

\subsubsection{The sign-flip mechanism.}

The key insight is that~$f$ need not be monotone in the same direction
for all domains.  Hajij et al.~\cite{hajij2025copresheaf} identify an
analogous phenomenon in graph neural networks: oversmoothing --- the
collapse of node representations after many message-passing layers ---
is the neural-network manifestation of \emph{over-consensus} in
multi-agent systems.  Their copresheaf networks resolve this by
relaxing the consistency constraints that sheaf-based architectures
impose.  In our framework, the signed laxator achieves the same
resolution algebraically: when~$\phi_G > 0$, the lax condition relaxes
sheaf consistency toward copresheaf-like diversity generation; when
$\phi_G < 0$, it tightens consistency toward consensus.

There are two competing effects of non-commutativity
on diversity:

\begin{enumerate}
\item \textbf{Constructive interference:} Non-commuting generators
  produce a richer orbit structure.  Agents connected by
  non-commuting channels explore more of the search space, sustaining
  diversity.  This is the dominant effect on \emph{rugged} landscapes
  (NK2--6, sudoku) where the fitness landscape has many local optima
  and the population benefits from sustained exploration.

\item \textbf{Destructive interference:} Non-commuting generators
  produce conflicting migration signals.  Agents receiving
  contradictory information from adjacent channels experience
  ``interference noise'' that disrupts coherent search.  This is the
  dominant effect on \emph{smooth} landscapes (OneMax) and tasks where
  the population maintains high diversity independently of topology
  (maze, where diversity $\approx 0.989$ regardless of arrangement).
  Sloboda~\cite{sloboda2026composition} (Thm.~6) gives this a precise
  algebraic cost: composing two sheaf morphisms, each of complexity
  $O(w)$, incurs an $\Omega(w^2)$ blowup --- meaning that the
  ``interference noise'' accumulates \emph{quadratically} in the depth
  of the non-commuting composition chain, not merely additively.
  Conversely, Hernandez Caralt et al.~\cite{hernandez2026identity} show
  that when $H^1 = 0$ (trivial cohomology), identity sheaves suffice
  and no such interference arises --- confirming that the quadratic
  cost is genuinely a feature of the non-trivial-cohomology regime.
  Yokoyama~\cite{yokoyama2026relative} further refines this threshold
  via \emph{relative} sheaf cohomology: even when $H^1 = 0$, the
  specification morphism can introduce relative obstructions
  $H^1_{\mathrm{rel}}(\mathcal{F}, W) \neq 0$, separating intrinsic
  topology obstructions from task-grounding obstructions.
  Hanks et al.~\cite{hanks2025relative}, implementing the program of
  Schmid~\cite{schmid2025prospectus}, give this relative approach
  empirical teeth: they show that $H^0_{\mathrm{rel}} = 0$ serves as a
  \emph{stability condition} for multi-agent systems, providing a
  computable certificate that no coordination failure can arise from the
  specification alone.  This complements Yokoyama's $H^1_{\mathrm{rel}}$
  obstruction: the relative cohomology ladder
  $H^0_{\mathrm{rel}} \to H^1 \to H^1_{\mathrm{rel}}$
  captures a progression from global stability to local inconsistency to
  task-grounded obstruction.

  At the next cohomological degree, Ghrist and
  Riess~\cite{ghrist2026cup} show that cup products in~$H^2$ measure
  \emph{interference between independent $H^1$ obstructions} --- that
  is, $H^2$ detects when two individually resolvable inconsistencies
  become jointly irresolvable.  They establish a paradox hierarchy
  $H^0$ (ambiguity) $\to$ $H^1$ (inconsistency) $\to$ $H^2$
  (inaccessibility), with a discrete Stokes theorem connecting
  boundary inconsistency to interior obstruction.  For the
  signed-laxator framework, this hierarchy is directly relevant:
  the constructive/destructive interference dichotomy of this section
  operates at the $H^1$ level (inconsistency of migration signals),
  but the \emph{interaction} between independent interference channels
  --- precisely what non-commutativity of Adjacent generators
  produces --- is an $H^2$ phenomenon.  The cup product structure
  may thus provide the algebraic foundation for the quadratic cost
  of Sloboda's Theorem~6.
\end{enumerate}

\noindent
On sudoku ($\eta^2 = 0.40$ for the arrangement effect at gen~500),
the constructive channel dominates: Separated generators, which
commute perfectly, allow additive diversity contributions without
cancellation, yielding the highest saturation ceiling.  Adjacent
generators, maximally non-commuting, produce the strongest
interference, \emph{lowering} the diversity ceiling through
destructive interaction of migration signals.  Hence
$\text{Separated} > \text{Adjacent}$ for the saturation metric.

On maze tasks, the population is already near-maximally diverse
($\approx 0.989$) and the topology effect is negligible
($\eta^2 \approx 0.01$).  In this regime, the small residual effect
of non-commutativity is to \emph{prevent} the population from
converging to a common solution -- the destructive interference of
Adjacent generators creates a weak diversity-\emph{preserving} force.
Hence the ordering reverses: $\text{Adjacent} > \text{Separated}$.

\subsubsection{Connection to Schur-convexity.}

This domain-dependent activation connects naturally to the
Schur-convexity framework of Amir, Bettini, and
Prorok~\cite{amir2026schur}.  Their key result: a multi-agent task
benefits from diversity (Schur-concave optimal) when the global
objective aggregates agent contributions via a Schur-concave function,
and benefits from homogeneity (Schur-convex optimal) when aggregation
is Schur-convex.

We identify the \emph{task curvature}
$\sigma \in [-1, +1]$ as the degree of Schur-concavity ($\sigma > 0$)
or Schur-convexity ($\sigma < 0$) of the task's aggregation function.
The connection to irrep activation is:

\begin{itemize}
\item \textbf{High $\sigma$ (Schur-concave, e.g., sudoku):}
  The task rewards diverse solutions.  The population dynamics project
  strongly onto irreps where the arrangement ordering matches the
  constructive-interference prediction: high commutator rank
  \emph{reduces} effective diversity by creating conflicting signals.
  Separated arrangements, being fully commutative, allow maximal
  constructive superposition.

\item \textbf{Low $\sigma$ (near-neutral, e.g., maze):}
  The task neither rewards nor penalizes diversity.  The population
  maintains near-maximal diversity regardless.  The residual
  arrangement effect is dominated by the destructive-interference
  channel: high commutator rank creates a weak force resisting
  convergence.  The active irreps are those where non-commutativity
  \emph{preserves} rather than \emph{reduces} diversity.

\item \textbf{Negative $\sigma$ (Schur-convex, e.g., OneMax):}
  The task rewards homogeneity.  The population dynamics select for
  convergence, and the arrangement effect operates through the
  interference channel.  Separated arrangements allow clean
  convergence; Adjacent ones create noisy signals that slow it.
  The partial match $\text{Separated} > \text{Nested} >
  \text{Adjacent}$ is consistent with this picture.
\end{itemize}

\subsubsection{The three-axis prediction.}

This analysis completes a three-axis framework for topology-mediated
optimization:

\begin{enumerate}
\item $\beta_1$ (first Betti number): \emph{how many} independent
  cycles exist in the communication topology.  Determines the
  \emph{capacity} for diversity maintenance.

\item $\lambda_2$ (Fiedler eigenvalue): \emph{how fast} information
  propagates through the topology.  Determines the \emph{timescale}
  of convergence.

\item $\sigma$ (task curvature / Schur-convexity): \emph{whether}
  diversity helps or hurts on the given task.  Determines which
  irreducible representation of the cactus group is active, and
  thereby the \emph{sign} of the arrangement effect.
\end{enumerate}

\noindent
The prediction chain is:
\[
  \sigma \;\longrightarrow\;
  w_\lambda(\mathcal{D})
  \;\longrightarrow\;
  \beta_1^*(\sigma)
  \;\longrightarrow\;
  \lambda_2^*(\beta_1),
\]
where $\beta_1^*(\sigma)$ is the optimal cycle rank for a task of
curvature~$\sigma$, and $\lambda_2^*(\beta_1)$ is the optimal
spectral gap for a graph with that cycle rank.

\begin{remark}[Composition complexity and the $\beta_1$ axis]
\label{rem:composition-complexity}
Sloboda~\cite{sloboda2026composition} (Thm.~6) establishes that the
composition of two sheaf morphisms, each of complexity $O(w)$,
produces an $\Omega(w^2)$ blowup.  This quantifies why $\beta_1$
matters structurally, not just combinatorially.  In a star topology
($\beta_1 = 0$), every coordination path passes through the hub at
length~$1$: there is no composition chaining and no quadratic cost.
In a cycle, agents compose migration signals \emph{sequentially}
around the ring, accumulating complexity at each step.  The
Fiedler eigenvalue~$\lambda_2$ controls the rate at which
inconsistency leaks across the cycle: a larger spectral gap forces
faster mixing, limiting how deeply the quadratic overhead can
accumulate before convergence.  This connects to the NK data:
NK2 shows a sustained topology advantage ($d = 0.56$) while NK4
is transient --- moderate ruggedness gives the landscape enough
structure that composition overhead matters before the population
saturates, but high ruggedness diversifies so rapidly that even
the $\Omega(w^2)$ cost is dominated by fitness-driven exploration.
\end{remark}

\begin{remark}[External validation: identity sheaves and trivial cohomology]
\label{rem:identity-sheaves}
Hernandez Caralt et al.~\cite{hernandez2026identity} provide independent
empirical validation of the signed-laxator thesis from the graph neural
network literature.  They show that \emph{identity sheaves} --- i.e.,
sheaves carrying no learned structure --- suffice for node classification
whenever the first cohomology vanishes ($H^1 = 0$), and identify a
``good heterophily $> 0.54$'' threshold as an empirical proxy for this
condition.  Read contrapositively, their result says: when $H^1 \neq 0$
(non-trivial cohomology), identity sheaves do \emph{not} suffice, and
the full sheaf machinery --- including the signed laxator --- becomes
necessary.  This is precisely the regime where our three-axis framework
predicts topology matters: non-trivial $\beta_1$ creates non-trivial
$H^1$, which activates the cactus-group interference structure and makes
the arrangement ordering observable.  The good-heterophily threshold can
be understood as marking the boundary between the commutative regime
(where separated/identity structure suffices) and the non-commutative
regime (where the full laxator is required to manage interference).
\end{remark}

\begin{remark}[Copresheaf duality and the laxator program]
\label{rem:copresheaf-duality}
Hajij et al.~\cite{hajij2025copresheaf} establish a formal duality
between sheaf neural networks (which enforce consistency via
restriction maps) and copresheaf neural networks (which propagate
information via extension maps without consistency constraints).
Their central finding --- that copresheaf architectures avoid the
oversmoothing pathology of deep sheaf networks --- provides a
structural parallel to the signed-laxator framework.  In our
setting, the laxator~$\phi_G$ can be reinterpreted as a bridge
between the sheaf and copresheaf perspectives on the communication
graph: the lax condition $\phi_{H} \circ F(f) \neq G(f) \circ
\phi_{G}$ is precisely the failure of strict sheaf consistency, and
the \emph{sign} of the laxator cell determines whether this failure
drives the system toward copresheaf-like exploration ($\phi_G > 0$)
or deeper sheaf-like consensus ($\phi_G < 0$).  The common receiving category construction of
Section~\ref{sec:copresheaf-future} (C.~Vega, private communication)
provides a concrete categorical foundation for this bridge, realizing
$\phi_G$ as an involution on the ambient section space.
\end{remark}

\subsubsection{Why the saturation formula fails.}

This framework explains the failure of the universal saturation
formula $\eta^2 \approx 1 - \exp(-\beta_1 / L)$ (where $L$ is
a landscape-dependent length scale), which achieves $R^2 = -0.75$
across domains~\cite{lyra2026topology}.  The formula implicitly
assumes a single, domain-independent relationship between $\beta_1$
and $\eta^2$.  But because different domains activate different
irreps, there is no such universal relationship: the effective weight
of $\beta_1$ depends on~$\sigma$, which varies by domain.

Concretely, on sudoku ($\sigma \gg 0$), $\beta_1$ has a strong
effect ($\eta^2 = 0.69$), while on maze ($\sigma \approx 0$),
$\beta_1$ has almost no effect ($\eta^2 = 0.02$), despite the
same graphs being used.  The formula correctly captures the
shape within a single domain (logistic $R^2 = 0.92$ for sudoku
alone) but fails catastrophically when pooled across domains because
the irrep weights~$w_\lambda$ change.

\subsubsection{Testable prediction.}

The irrep-activation hypothesis generates a concrete, falsifiable
prediction:

\begin{prediction}
\label{pred:per-irrep-reversal}
Compute the per-irrep commutator rank ordering separately for each
non-trivial $V_\lambda$ of $S_n$.  For irreps activated by
Schur-concave tasks (sudoku), the \emph{observable} arrangement
effect should follow the standard ordering
$\mathrm{Separated} > \mathrm{Adjacent}$ (because commutativity
enables constructive superposition).  For irreps activated by
near-neutral tasks (maze), the observable arrangement effect should
reverse to $\mathrm{Adjacent} > \mathrm{Separated}$ (because
non-commutativity preserves diversity).
\end{prediction}

\noindent
Testing this prediction requires two steps: (a)~identify the
domain-dependent weights~$w_\lambda(\mathcal{D})$ from the population
Fourier spectrum on each task, and (b)~verify that the weighted
commutator rank, computed via
equation~\eqref{eq:weighted-arrangement}, reproduces the observed
arrangement ordering on each domain.  Step~(a) requires
Fourier analysis on~$S_n$ of the population diversity vectors, which
is computationally feasible for the group sizes in our experiments
($n = 4, 5, 8$).

\subsubsection{The common receiving category and the signed involution.}
\label{sec:copresheaf-future}

The sheaf--copresheaf duality of Hajij et
al.~\cite{hajij2025copresheaf} suggests a natural next step for the
signed-laxator program: formalize~$\phi_G$ as a morphism in a
category that carries both sheaf and copresheaf structure
simultaneously.  The construction below, due to C.~Vega (private
communication), provides a concrete realization.

\paragraph{The common receiving category.}
Let $G = (V,E)$ be a communication graph and $\mathbb{F}$ a field
(we work over $\mathbb{Q}$ throughout).  A sheaf
$\mathcal{F} \in \mathrm{Sh}(G;\mathbb{F})$ assigns an
$\mathbb{F}$-vector space~$\mathcal{F}(v)$ to each node~$v$ with
restriction maps on edges enforcing consistency, while a copresheaf
$\mathcal{G} \in \mathrm{CoSh}(G;\mathbb{F})$ assigns vector spaces
with extension maps enabling directed flow without consistency
constraints~\cite{hajij2025copresheaf}.  Both embed into a common
target: the category of \emph{graded sections}
\[
  \Gamma_{\mathrm{gr}}(G;\mathbb{F})
  \;=\;
  \bigoplus_{v \in V} \mathcal{F}(v),
\]
the underlying graded vector space obtained via a forgetful functor
that discards all compatibility data.  The sheaf functor lands in the
\emph{equalizer image} --- the subspace of consistent sections
satisfying all restriction conditions --- while the copresheaf functor
lands in the \emph{colimit image} --- the quotient carrying
forward-propagated sections.  The laxator $\phi_G$ is the comparison
map between these two sub-objects of $\Gamma_{\mathrm{gr}}$.

\paragraph{The signed involution.}
Let $P_+$ denote the orthogonal projection onto the equalizer image
(the sheaf-consistent subspace) and $P_-$ the projection onto its
complement in $\Gamma_{\mathrm{gr}}$ (the colimit complement).  When
$P_+$ and $P_-$ are orthogonal idempotents --- that is,
$P_+^2 = P_+$, $P_-^2 = P_-$, $P_+ P_- = P_- P_+ = 0$, and
$P_+ + P_- = I$ --- the comparison map
\begin{equation}
\label{eq:phi-involution}
  \phi_G \;=\; P_+ \;-\; P_-
\end{equation}
satisfies $\phi_G^2 = \mathrm{id}$.  This decomposition splits the
section space into a $+1$ eigenspace (sheaf-like: global sections,
$H^0$) and a $-1$ eigenspace (copresheaf-like: obstructed modes
that violate consistency).  The sign of $\phi_G$ on a given section
records whether that section is driven toward consensus or toward
divergence.

\paragraph{Computational verification on $P_3$.}
We verify the involution property on the path graph
$P_3 = (v_1 - v_2 - v_3)$ with rank-$2$ fibers
($\mathcal{F}(v_i) = \mathbb{Q}^2$ for each~$i$), the simplest
non-trivial case.  All computations are exact over~$\mathbb{Q}$.
\begin{enumerate}
\item $\phi_G^2 = I_{6 \times 6}$, confirming the involution
  property~\eqref{eq:phi-involution} exactly.
\item $P_+$ and $P_-$ are orthogonal idempotents:
  $P_+^2 = P_+$, $P_-^2 = P_-$, $P_+ P_- = P_- P_+ = 0$.
\item Every $2 \times 2$ block of $\phi_G$ is a scalar multiple of
  $I_2$, revealing complete decoupling of the rank-$2$ fibers.
\item The $+1$ eigenspace has dimension~$2$, corresponding to
  $H^0(P_3;\mathcal{F}) = \mathbb{Q}^2$; the $-1$ eigenspace has
  dimension~$4$.
\item $H^1(P_3;\mathcal{F}) = 0$, as expected for a tree graph.
\end{enumerate}
The $P_3$ result is initial verification on the simplest possible
graph --- a tree with trivial first cohomology.  The test extends
naturally to graphs with non-trivial~$H^1$ (e.g., the cycle $C_n$),
where the $-1$ eigenspace should acquire additional structure reflecting
the cohomological obstruction.

\paragraph{Connection to Hecke structure.}
The orthogonal idempotent decomposition $\phi_G = P_+ - P_-$ is
formally identical to a structure in the Hecke algebra
(C.~Vega, private communication).  In the Hecke algebra
$\mathcal{H}_q(S_n)$, the symmetrizer decomposes as
$\Pi_q = (1+q)^N \cdot \prod_i e_i$, where each $e_i^2 = e_i$.
The $1/3$ uniform averaging that appears in the $P_3$ computation
--- each vertex contributing equally to the global section ---
recalls the group algebra idempotent $(1/|G|) \cdot \sum_g g$.
This parallel suggests that the signed involution may extend to a
$q$-deformation, with the Hecke parameter~$q$ playing the role of
the task curvature~$\sigma$ from Section~\ref{sec:domain-dependent-activation}:
$q = 1$ recovers the symmetric (sheaf-like) regime, $q = 0$
the degenerate (copresheaf-like) regime.

\paragraph{The vanishing condition.}
When $H^1(G;\mathcal{F}) = 0$ --- equivalently, when identity sheaves
suffice in the sense of Hernandez Caralt et
al.~\cite{hernandez2026identity} --- the equalizer image spans all of
$\Gamma_{\mathrm{gr}}$: every section is consistent, $P_+ = I$,
$P_- = 0$, and $\phi_G = I$.  The comparison map is an isomorphism
and no laxator is needed.  This confirms that the signed-laxator
framework is activated precisely in the non-trivial cohomology regime.
Conversely, when $H^1 \neq 0$, the $-1$ eigenspace is non-trivial,
and the sign of $\phi_G$ on a given section measures the orientation
mismatch between the sheaf and copresheaf images: positive when they
agree (the section respects consistency), negative when reversed (the
section violates it).  The three sign-flip observations of Section~3
can thus be understood as manifestations of different sections falling
on opposite sides of this eigenspace
decomposition~\cite{hajij2025copresheaf,hernandez2026identity,sloboda2026composition,ghrist2026cup,hanks2025relative}.

\paragraph{Algorithmic resolution via holonomy.}
The HATCC algorithm of ter~Horst, Mahadevan, and
Zambrano~\cite{terhorst2026categorical} provides a natural algorithmic
complement to our obstruction theory.  Where our signed restriction maps
characterise \emph{when} composition degrades diversity ($H^1 \neq 0$),
HATCC provides a deterministic procedure for \emph{detecting} such
obstructions via holonomy computation and \emph{resolving} them by
compiling non-trivial holonomy into mode variables that partition the
problem into exact sub-problems.  Their key result --- that the augmented
graph after HATCC compilation is always a tree ($\beta_1 = 0$) --- means
the obstruction is resolved structurally rather than iteratively.
Adapting HATCC to signed sheaves, where transport kernels must respect
the sign structure $\sigma = \pm 1$ of restriction maps, is a natural
direction for future work.

% ================================================================
% NOTE FOR CLIO:
%
% Remaining gaps after Section 5.2.8 rewrite:
%
% 1. Section 5.3 (fiber rotation mechanism) should connect the
%    domain-dependent sign to the monodromy of the flat connection
%    as kappa varies.  Does the sign-flip correspond to the fiber
%    rotating through a singular point on the Riemann sphere?
%
% 2. The Adjacent commutator rank being *exactly* constant (no
%    variance at all within each irrep) feels like it should have
%    a clean proof.  Is this related to the palindromic structure
%    you mentioned in the eigenvalue-positivity work?
%
% 3. Extend the P_3 verification to C_4 or C_5 (non-trivial H^1).
%    The -1 eigenspace should pick up cohomological structure.
% ================================================================
